\subsection{Synthetic Experiments}

We first carry out experiments on several synthetic text datasets consisting of five Tomita grammars of binary strings and a context-free ``Motzkin'' grammar over the trinary alphabet $\Sigma_{M} = \{\mathtt{\ (\ ,\ \starchar\ ,\ )\ }\}$~\citep{tomita1982, alexander2018}. The latter consists of all strings whose parentheses are properly balanced, with no constraints placed on the $\mathtt{\starchar}$ characters. 

In each case we train the u-MPS on strings of a restricted length from the grammar and then sample new strings of unseen lengths from the trained u-MPS, with the model assessed on the percentage of sampled strings which match the grammar. The sampling comes in two forms, either fixed length-$n$ sampling (corresponding to $R = \Sigma^n$), or character completion sampling, where a single character in a reference string is masked and the prefix and suffix $p$ and $s$ are used to guess it (corresponding to $R = p \Sigma s$). While more general sampling experiments can easily be imagined, we have chosen these tasks because they allow for direct comparisons with a variety of baselines, including (unidirectional and bidirectional) LSTMs, HMMs, and Transformers.

While unbiased fixed-length sampling is easy for u-MPS via Algorithm~\ref{alg:sampling}, we found that the unidirectional LSTM baseline required an additional positional encoding in its inputs to avoid rapid degeneration in the output text when sampling past the longest length seen in training. At sampling time, we vary the length scale associated with this encoding based on the desired sampling length, so that the final step of sampling is always associated with the same positional encoding vector.

We train the u-MPS, LSTM, HMM, and small Transformer models using gradient descent on a negative log likelihood (NLL) loss with the Adam~\citep{kingma2015} optimizer. For each experiment we use models of $D=20$ and $D=50$ hidden units, with LSTMs chosen to have one layer and Transformers with two layers and 4 heads. Five independent trials are used for each value of $D$, with the final validation loss used to select the best model for generating samples. We use a piecewise constant learning rate between $10^{-2}$ and $10^{-4}$, and early stopping to choose the end of training.

In the Tomita experiments~(Table~\ref{tab:tomita}) u-MPS give impressive performance, in many cases achieving perfect accuracy in sampling strings of unseen sizes within the language. This is true not only in the simpler grammars Tomita 3 and 4, but also in the more difficult Tomita 5, where valid strings satisfy the nonlocal constraint of containing an even number of \texttt{0}'s and of \texttt{1}'s. The HMM and LSTM also attain reasonably high sample accuracy, although in a manner that degrades faster with sequence length than the u-MPS, while the Transformer performs worst.

Similar results are seen with the context-free Motzkin language~(Table~\ref{tab:motzkin}), where a fixed-length sampling task similar to the Tomita experiment is paired with a character completion task. A separate bidirectional LSTM must be used for this latter task, since unidirectional LSTMs cannot make use of future context information. By contrast, a trained u-MPS model can be employed in both of these settings without any task-specific adaptation, as well as in more general sentence completion tasks involving connected or disjoint regions of missing text~(tasks which cannot be easily handled by standard RNN models). The u-MPS does substantially better in generalizing and reproducing the structure of Motzkin strings than the unidirectional LSTM, HMM, and Transformer, being able to sample strings of over 3 times the length seen in training with over 90\% accuracy. The u-MPS is outperformed only by the bidirectional LSTM in character completion experiments on smaller training sets.

Given the ability of HMMs to exactly reproduce the distributions associated with Tomita languages and (bounded length) Motzkin languages, it is surprising that the u-MPS still manages to more easily learn such distributions in practice. Somewhat surprising also is the poor performance of the Transformer models, which is likely a result of the small sizes of the string datasets used.

We additionally benchmark the relative runtime of sequential and parallel evaluation for u-MPS running on a CPU or GPU (Table~\ref{tab:runtime}). For a u-MPS of $D$ hidden states with strings of length $n$, RNN-style sequential evaluation has parallel depth $\mathcal{O}(n)$ and cost $\mathcal{O}(nD^2)$, while parallel evaluation trades this for a parallel depth of $\mathcal{O}(\log n)$ and cost $\mathcal{O}(nD^3)$. This would suggest parallel evaluation having benefits for the total runtime when hardware acceleration is present, which the runtimes in Table~\ref{tab:runtime} confirm.

